% Осень26!B15; нумерация внутри столбца без пустых строк.
% Источник: лекция 2, §6 (строка 175)
\definitionticket{12}{Классический объем в \(\mathbb R^m\)}
\label{definition:12}

\subsubsection*{\texorpdfstring{Классический объём в \(\mathbb{R}^m\)}{Классический объём в \textbackslash mathbb\{R\}\^{}m}}

В \(\mathbb{R}^m\) --- полукольцо ячеек \(\mathcal{P}^m\) (полуоткрытые параллелепипеды). \textbf{Классический объём}: \[\mu[a,b)=\prod_{i=1}^{m}(b_i-a_i)\] --- произведение длин проекций (рёбер) параллелепипеда, «не обращая внимания на граничные эффекты»: часть границы взяли, часть нет --- на объём это не влияет, граница «в разумном смысле окажется множеством нулевого объёма».

Лектор аддитивность \textbf{не доказывал} («какие-то фокусы, идейно ненасыщенные») и дал идею.

\textbf{Идея доказательства аддитивности.} Достаточно понять случай разбиения параллелепипеда сеткой: если по одной оси кусочки имеют длины \(h_i\), по другой \(v_j\), то раскрытие скобок в произведении сумм \[(\textstyle\sum_i h_i)\bigl(\textstyle\sum_j v_j\bigr)=\sum_{i,j}h_iv_j\] перечисляет в точности все клеточки разбиения --- сумма объёмов кусочков равна объёму целого. «Глядя на тетрадку в клеточку, все верят, что это адекватная аддитивная функция». Общий случай сводится к подразбиениям: при подразбиении параллелепипеда сумма объёмов сохраняется, а любое разбиение дополняется до сеточного.
